\subsection{Mailboxes and Mailbox Patterns}
\label{sec:messages}

Like Pat but unlike in the original definition of mailbox types
\cite{deliguoroMailboxTypesUnordered2018}, we do not include message payloads as
part of a mailbox type, but rely only on \emph{message tags}
\msg{m} like \msg{Put} and \msg{Get} and associate types with messages when defining the typing rules for
Pat (\S\ref{sec:typing-rules}).

% that should be possible with some effort.

% To simplify reasoning about messages, we restrict the set of possible messages to a predefined but arbitrary set.
% Additionally, equality of messages must be decidable.

% \begin{linkdef}{https://edgardschi.github.io/mailbox-types-rocq/MailboxTypes.Message.html\#IMessage}[Message]
% \todo{ES}{Remove this definition, just mention in text}
% A message $\msg{m}, \msg{n}$ is a value of an arbitrary type equipped with decidable equality.
% \end{linkdef}

\begin{figure}[t]
  \centering
  \begin{minipage}{0.7\textwidth}
    \begin{tabular}{lccl}
      Mailbox configurations (\rocqlink{https://edgardschi.github.io/mailbox-types-rocq/MailboxTypes.Message.html\#MessageConfig_def}) \; & $m$ \; & $::=$ \; & $\mbempty \mid \mbconf{\msg{m}} \mid m_1 \mbu m_2$\\
      Mailbox patterns (\rocqlink{https://edgardschi.github.io/mailbox-types-rocq/MailboxTypes.MailboxPatterns.html\#MPattern}) \; & $E, F$ \; & $::=$ \; & $\mbzero \mid \mbone \mid \mbmsg{m} \mid E \mbcomp F$\\
       \; & \; & $\mid$ & $E \mbchoice F \mid \mbstar{E}$
    \end{tabular}
  \end{minipage}
  \hfill
  \begin{minipage}{0.27\textwidth}
    Pattern power (\rocqlink{https://edgardschi.github.io/mailbox-types-rocq/MailboxTypes.MailboxPatterns.html\#Pow})\\
    \begin{tabular}{lcl}
      $E^0$ & $:=$ & $\mbone$\\
      $E^{n+1}$ & $:=$ & $E \mbcomp E^n$
    \end{tabular}
  \end{minipage}

  \caption{The syntax of mailbox configurations and patterns.}
  \label{fig:syntax-mailbox-patterns}
\end{figure}

% In Rocq this definition can be achieved using a type class.
% \begin{coq}
% Class IMessage Message : Type :=
% {
%   eq_dec : forall m n : Message, {m = n} + {m <> n}
% }.
% \end{coq}
% This approach has several advantages over statically defining a concrete message type.
% Firstly, we can argue about any kind of messages and are not restricted by the initial definition of messages.
% Furthermore, we can define our own message type or use existing data types.
% An additional advantage of this approach is that needed properties about messages can be extended without much effort.
% For example, if we extend messages to carry more than one payload, we could include a function mapping each message to the amount of data it carries:
% \begin{coq}
% Class IMessage Message : Type :=
% {
%   eq_dec : forall m n : Message, {m = n} + {m <> n};
%   content_size : Message -> $\mathbb{N}$
% }.
% \end{coq}
%
% As an example of how one can define messages, consider strings.
% We can use Rocq's predefined implementation of strings and use them as messages by proving the corresponding type class instance:
% \begin{coq}
% Instance StringMessage : IMessage string :=
% {
%   eq_dec := string_dec
% }.
% \end{coq}
%
% As an example of custom defined messages, consider the following type definition that represents two messages:
% \begin{coq}
% Inductive SendReceive : Type :=
%     Send : SendReceive
%   | Receive : SendReceive.
% \end{coq}
% By proving that equality is decidable for this given type, we can define the message type class instance.
% \begin{coq}
% Instance SendReceiveMessage : IMessage SendReceive :=
% {
%   eq_dec := SendReceive_eq_dec
% }.
% \end{coq}

\paragraph{Mailbox Configurations.}
A mailbox represents an unordered collection of messages, where a message can be included multiple times.
Figure \ref{fig:syntax-mailbox-patterns} shows the definition of a \emph{mailbox configuration} $m$.
It is either the empty mailbox $\mbempty$, a singleton mailbox $\mbconf{\msg{m}}$ containing only message $\msg{m}$, or the union of two mailboxes $m_1 \mbu m_2$.
In Rocq, we model mailbox configurations as lists, where $\mbu$ is implemented as append.
To reason about the equality of mailboxes, i.e.\ mailboxes with the same content, we need to take into account that order does not matter.
% For example, mailboxes $\mbconf{\msg{m}} \mbu \mbconf{\msg{m}} \mbu \mbconf{\msg{n}}$ and $\mbconf{\msg{m}} \mbu \mbconf{\msg{n}} \mbu \mbconf{\msg{m}}$
% should be considered equivalent.
Mailbox configurations $m_1$ and $m_2$ are equal, written $m_1 \approxident m_2$, if they are permutations of each other (\rocqlink{https://edgardschi.github.io/mailbox-types-rocq/MailboxTypes.Message.html\#mailbox_eq}).
% We define the relation $m_1 \approxident m_2$ to be holding, if $m_1$ and $m_2$ are permutations of each other.
We then have $\mbconf{\msg{m}} \mbu \mbconf{\msg{m}} \mbu \mbconf{\msg{n}} \approxident \mbconf{\msg{m}} \mbu \mbconf{\msg{n}} \mbu \mbconf{\msg{m}}$,
but $\mbconf{\msg{n}} \mbu \mbconf{\msg{n}} \not\approxident \mbconf{\msg{n}}$.


\paragraph{Mailbox Patterns.}
The crucial building block of mailbox types are \emph{mailbox patterns}.
They can be thought of as a commutative regular expression that describes the content of a mailbox.
Figure \ref{fig:syntax-mailbox-patterns} shows the definition of mailbox patterns.
The pattern $\mbzero$ denotes an \emph{unreliable} mailbox, i.e.\ one that
received an unexpected message and which cannot be used for further
communication.  Pat's type system ensures that it is not possible to create such
a mailbox using well-typed terms.
%
Pattern $\mbone$ represents an \emph{empty} mailbox, containing zero messages.
%
The pattern $\mbmsg{m}$ denotes a mailbox containing a \emph{single message} $\msg{m}$.
Mailbox \emph{composition} is described by pattern $E \mbcomp F$, combining both patterns $E$ and $F$.
As an example, pattern $\mbmsg{m} \mbcomp \mbmsg{n}$ would describe a mailbox containing both $\msg{m}$ and $\msg{n}$.
%
On the other hand, $E \mbchoice F$ denotes \emph{choice} between patterns.
In particular, it describes that the mailbox is either described by pattern $E$ or by pattern $F$.
Lastly, pattern $\mbstar{E}$ describes repeated composition of pattern $E$.
For example, $\mbstar{\mbmsg{m}}$ denotes a mailbox containing zero or more $\msg{m}$ messages.

\paragraph{Pattern Semantics.}
Since patterns describe the contents of an unordered mailbox, multiple patterns can describe the same mailbox.
For example, pattern $\mbmsg{m} \mbcomp \mbmsg{n}$ describes the same mailbox as $\mbmsg{n} \mbcomp \mbmsg{m}$.
Thus, checking equality of mailbox patterns requires a semantic definition.
The semantics of mailbox types are usually given by a denotational semantics
that map patterns to sets of multisets of atoms~\cite{deliguoroMailboxTypesUnordered2018}.
Since we are working with Rocq, which is based on constructive logic, we require
a method of constructing these sets;
a particular challenge is representing the $\mbstar{E}$ pattern whose denotation
is an infinite union of multisets and not easily representable in Rocq.
Instead, we use a relational approach to connect mailbox configurations to
patterns. 

\begin{figure}[t]
  % \textbf{Pattern semantics} \hfill \fbox{$m \in E$}
  \begin{center}
        \begin{prooftree}
            \infer0{\mbempty \in \mbone}
        \end{prooftree}
        \quad
        \begin{prooftree}
            \infer0{\mbconf{\msg{m}} \in \mbmsg{m}}
        \end{prooftree}
        \quad
        \begin{prooftree}
            \hypo{m \in E}
            \infer1{m \in E \mbchoice F}
        \end{prooftree}
        \quad
        \begin{prooftree}
            \hypo{m \in F}
            \infer1{m \in E \mbchoice F}
        \end{prooftree}
        \quad
        \begin{prooftree}
            \hypo{\exists n.\ m \in E^n}
            \infer1{m \in \mbstar{E}}
        \end{prooftree}
        \quad
        \begin{prooftree}
            \stackedhypos{1}{
              \hypo{m \mbeq m_1 \mbu m_2}
            }{2}{
              \hypo{m_1 \in E}
              \hypo{m_2 \in F}
            }
            \infer1{m \in E \mbcomp F}
        \end{prooftree}
    \end{center}

    \caption{The semantics of mailbox patterns (\rocqlink{https://edgardschi.github.io/mailbox-types-rocq/MailboxTypes.MailboxPatterns.html\#valueOf}).}
    \label{fig:pattern-semantics}
\end{figure}

Instead of building possibly infinite sets of possible configurations, we relate a mailbox configuration to a pattern, written $m \in E$.
A configuration is included in the set induced by a pattern if the configuration can be constructed using the pattern.
For example, we have $\mbconf{\msg{m}} \mbu \mbconf{\msg{n}} \in \mbmsg{m} \mbcomp \mbmsg{n}$ and $\mbconf{\msg{n}} \mbu \mbconf{\msg{m}} \in \mbmsg{m} \mbcomp \mbmsg{n}$,
but also $\mbconf{\msg{m}} \mbu \mbconf{\msg{n}} \in \mbmsg{m} \mbcomp \mbmsg{n} \mbcomp \mbone$.

The inference rules for this relation are given in Figure~\ref{fig:pattern-semantics}.
The first two rules describe that the empty mailbox configuration $\mbempty$ is the pattern $\mbone$ and a configuration containing a single message $\msg{m}$ is part of pattern $\mbmsg{m}$.
For the choice operator $E \mbchoice F$, there are two rules: one for when $m$ is part of $E$ and one for when $m$ is part of $F$.
The rule for composition is the most complex one.
If $m$ is part of a composition of patterns $E$ and $F$, then $m$ is required to be split into two sub-configurations, where one is part of $E$ and one is part of $F$.
Lastly, the rule for $\mbstar{E}$ uses the 
\emph{power} operation (defined in Figure~\ref{fig:syntax-mailbox-patterns})
to state that any mailbox configuration that can be built using a pattern $E^n$
(for some $n$) is included in $\mbstar{E}$.

One important observation to make is that we do not include a rule for pattern $\mbzero$.
According to the definition of de'Liguoro and Padovani \cite{deliguoroMailboxTypesUnordered2018}, $\mbzero$ should map to the empty set, containing no configurations.
Thus, since we are relating configurations and patterns, there must be no relation between any configuration and $\mbzero$.

We define \emph{pattern inclusion} $E \mbpinc F$ as the proposition $\forall m .\ m \in E \Rightarrow m \in F$ (\rocqlink{https://edgardschi.github.io/mailbox-types-rocq/MailboxTypes.MailboxPatterns.html\#MPInclusion}).
% \emph{Pattern inclusion} $E \mbpinc F$ is defined not in terms of subsets, but rather as a proposition $\forall m \in E \Rightarrow m \in F$.
Taking inclusion in both directions yields pattern \emph{equivalence} $E \mbpeq F$ (\rocqlink{https://edgardschi.github.io/mailbox-types-rocq/MailboxTypes.MailboxPatterns.html\#MPEqual}).

To verify this formalization has the required properties described by de'Liguoro and Padovani \cite{deliguoroMailboxTypesUnordered2018}, we prove
the following properties:

\begin{lemma}
  The relation $\mbpinc$ is a precongruence, i.e. the following properties hold:
  \begin{itemize}
    \item Reflexive (\rocqlink{https://edgardschi.github.io/mailbox-types-rocq/MailboxTypes.MailboxPatterns.html\#MPInclusion_refl}): $E \mbpinc E$
    \item Transitive (\rocqlink{https://edgardschi.github.io/mailbox-types-rocq/MailboxTypes.MailboxPatterns.html\#MPInclusion_trans}): if $E \mbpinc F$ and $F \mbpinc G$, then $E \mbpinc G$
    \item Compatible wrt. $\mbcomp$ (\rocqlink{https://edgardschi.github.io/mailbox-types-rocq/MailboxTypes.MailboxPatterns.html\#MPInclusion_compatable_Comp}): if $E_1 \mbpinc F_1$ and $E_2 \mbpinc F_2$, then $(E_1 \mbcomp E_2) \mbpinc (F_1 \mbcomp F_2)$
    \item Compatible wrt. $\mbchoice$ (\rocqlink{https://edgardschi.github.io/mailbox-types-rocq/MailboxTypes.MailboxPatterns.html\#MPInclusion_compatable_Choice}): if $E_1 \mbpinc F_1$ and $E_2 \mbpinc F_2$, then $(E_1 \mbchoice E_2) \mbpinc (F_1 \mbchoice F_2)$
  \end{itemize}
\end{lemma}

\begin{lemma}
  The relation $\mbpeq$ is an equivalence relation, i.e. it is
  \begin{itemize}
    \item Reflexive (\rocqlink{https://edgardschi.github.io/mailbox-types-rocq/MailboxTypes.MailboxPatterns.html\#MPEqual_refl}): $E \mbpeq E$
    \item Transitive (\rocqlink{https://edgardschi.github.io/mailbox-types-rocq/MailboxTypes.MailboxPatterns.html\#MPEqual_trans}): if $E \mbpeq F$ and $F \mbpeq G$, then $E \mbpeq G$
    \item Symmetric (\rocqlink{https://edgardschi.github.io/mailbox-types-rocq/MailboxTypes.MailboxPatterns.html\#MPEqual_sym}): if $E \mbpeq F$, then $F \mbpeq E$
  \end{itemize}
\end{lemma}

Additionally, we show that we indeed work with a commutative Kleene algebra $(M, \mbchoice, \mbcomp, \mbstar{}, \mbzero, \mbone)$,
where $M$ is the set of messages (\rocqlink{https://edgardschi.github.io/mailbox-types-rocq/MailboxTypes.MailboxPatterns.html\#MPattern_props}).%, as described by de'Liguoro and Padovani.

% \begin{lemma}
%   \todo{ES}{Remove this, just mention in text}
%   The relation $\mbpinc$ is a \emph{precongruence}, i.e.\ it is
%   \begin{itemize}
%     \rocqitem{https://edgardschi.github.io/mailbox-types-rocq/MailboxTypes.MailboxPatterns.html\#MPInclusion_refl} Reflexive: $E \mbpinc E$
%     \rocqitem{https://edgardschi.github.io/mailbox-types-rocq/MailboxTypes.MailboxPatterns.html\#MPInclusion_trans} Transitive: if $E \mbpinc F$ and $F \mbpinc G$, then $E \mbpinc G$
%     \rocqitem{https://edgardschi.github.io/mailbox-types-rocq/MailboxTypes.MailboxPatterns.html\#MPInclusion_compatable_Comp}
%       Compatible with composition: if $E_1 \mbpinc F_1$ and $E_2 \mbpinc F_2$, then $(E_1 \mbcomp E_2) \mbpinc (F_1 \mbcomp F_2)$
%     \rocqitem{https://edgardschi.github.io/mailbox-types-rocq/MailboxTypes.MailboxPatterns.html\#MPInclusion_compatable_Choice}
%       Compatible with composition: if $E_1 \mbpinc F_1$ and $E_2 \mbpinc F_2$, then $(E_1 \mbchoice E_2) \mbpinc (F_1 \mbchoice F_2)$
%   \end{itemize}
% \end{lemma}
%
% \begin{lemma}
%   \todo{ES}{Remove this, just mention in text}
%   The relation $\mbpeq$ is an \emph{equivalence relation}, i.e.\ it is
%   \begin{itemize}
%     \rocqitem{https://edgardschi.github.io/mailbox-types-rocq/MailboxTypes.MailboxPatterns.html\#MPEqual_refl} Reflexive: $E \mbpeq E$
%     \rocqitem{https://edgardschi.github.io/mailbox-types-rocq/MailboxTypes.MailboxPatterns.html\#MPEqual_sym} Symmetric: if $E \mbpeq F$, then $F \mbpeq E$
%     \rocqitem{https://edgardschi.github.io/mailbox-types-rocq/MailboxTypes.MailboxPatterns.html\#MPEqual_trans} Transitive: if $E \mbpeq F$ and $F \mbpeq G$, then $E \mbpeq G$
%   \end{itemize}
% \end{lemma}

% \begin{figure}
%   \begin{minipage}{.49\linewidth}
%   \begin{align*}
%     E \mbchoice \mbzero &\mbpeq E\\
%     E \mbchoice E &\mbpeq E\\
%     E \mbchoice F &\mbpeq F \mbchoice E\\
%     E \mbchoice (F \mbchoice F') &\mbpeq (E \mbchoice F) \mbchoice F'\\
%   \end{align*}
%   \end{minipage}
%   \begin{minipage}{.49\linewidth}
%   \begin{align*}
%     E \mbcomp \mbone &\mbpeq E\\
%     E \mbcomp \mbzero &\mbpeq \mbzero\\
%     E \mbcomp F &\mbpeq F \mbcomp E\\
%     E \mbcomp (F \mbcomp F') &\mbpeq (E \mbcomp F) \mbcomp F'\\
%   \end{align*}
%   \end{minipage}
%   \vspace{-.5cm}
%   \begin{align*}
%     E \mbcomp (F \mbchoice F') &\mbpeq (E \mbcomp F) \mbchoice (E \mbcomp F')\\
%     \mbstar{E} &\mbpeq \mbstar{E} \mbchoice E \mbcomp \mbstar{E}
%   \end{align*}
%   \caption{\rocqlink{https://edgardschi.github.io/mailbox-types-rocq/MailboxTypes.MailboxPatterns.html\#MPattern_props}{Properties of mailbox patterns}}
%   \label{fig:pattern-properties}
% \end{figure}

% We are now ready to prove that we indeed work with a commutative Kleene algebra, as described by de'Liguoro and Padovani.
% Figure \ref{fig:pattern-properties} shows the necessary properties required.
% We see that both $\mbchoice$ and $\mbcomp$ are commutative and associative.
% Additionally, both a have unit elements and $\mbcomp$ distributes over $\mbchoice$.
% We also can observe the recursive nature of $\mbstar{}$, which can be unfolded an arbitrary amount of times.

\begin{figure}[t]
    {\small
  \textbf{Pattern residuals} (\rocqlink{https://edgardschi.github.io/mailbox-types-rocq/MailboxTypes.MailboxPatterns.html\#PatternResidual}) \hfill \fbox{$\mbpres{E}{m}{F}$}
  \begin{center}
        \begin{prooftree}
          \infer0{\mbpres{\mbzero}{m}{\mbzero}}
        \end{prooftree}
        \qquad
        \begin{prooftree}
          \infer0{\mbpres{\mbone}{m}{\mbzero}}
        \end{prooftree}
        \qquad
        \begin{prooftree}
          \infer0{\mbpres{\mbmsg{m}}{m}{\mbone}}
        \end{prooftree}
        \qquad
        \begin{prooftree}
          \hypo{\msg{m} \neq \msg{n}}
          \infer1{\mbpres{\mbmsg{m}}{n}{\mbzero}}
        \end{prooftree}
        \qquad
        \begin{prooftree}
            \hypo{\mbpres{E}{m}{E'}}
            \infer1{\mbpres{\mbstar{E}}{m}{E' \mbcomp \mbstar{E}}}
        \end{prooftree}
    \end{center}
    \begin{center}
        \begin{prooftree}
            \hypo{\mbpres{E}{m}{E'}}
            \hypo{\mbpres{F}{m}{F'}}
            \infer2{\mbpres{E \mbchoice F}{m}{E' \mbchoice F'}}
        \end{prooftree}
        \qquad
        \begin{prooftree}
            \hypo{\mbpres{E}{m}{E'}}
            \hypo{\mbpres{F}{m}{F'}}
            \infer2{\mbpres{E \mbcomp F}{m}{(E' \mbcomp F) \mbchoice (E \mbcomp F')}}
        \end{prooftree}
    \end{center}
}
  \caption{The definition of pattern residuals.}
    \label{fig:pattern-residuals}
\end{figure}

\paragraph{Pattern Residuals.}
Mailboxes are dynamic objects, not static.
Their content changes dynamically while receiving and consuming messages.
For example, when consuming message \msg{m} from a mailbox described by pattern $\mbmsg{m}$ the resulting mailbox should be empty,
described by pattern $\mbone$.
To model this behavior, de'Liguoro and Padovani introduce the concept of \emph{pattern residuals} \cite{deliguoroMailboxTypesUnordered2018} (Figure \ref{fig:pattern-residuals}).
In Pat \cite{fowlerSpecialDeliveryProgramming2023a}, pattern residuals correspond exactly to the Brzozowski
derivative of a commutative regular expression \cite{brzozowskiDerivativesRegularExpressions1964}.
A pattern residual describes the resulting pattern when a message is consumed, i.e.\ removed from the mailbox.
% The original description of pattern residuals defines them as a partial operator on patterns \cite{deliguoroMailboxTypesUnordered2018}.
% \sjfnote{I think Pat's derivative operator is total, right, as it corresponds to
% the Brzozowski derivative?}
% Since Rocq does not allow partial functions in the traditional sense, we define a residual as a relation between a pattern, a message and the result pattern.

An important property of residuals is \emph{balancing}.
For a configuration described by pattern $\mbmsg{m} \mbcomp E$, which is included in a pattern $F$, $E$ should be included in the residual $\mbpres{F}{m}{F'}$.
This property is needed for reasoning about the type of a mailbox after a
message has been received.

\begin{lemma}[Balancing (\rocqlink{{https://edgardschi.github.io/mailbox-types-rocq/MailboxTypes.MailboxPatterns.html\#Balancing}})]
  If $\mbmsg{m} \mbcomp F \mbpinc E$ where $F \not\mbpinc \mbzero$ and $\mbpres{E}{m}{E'}$, then $F \mbpinc E'$.
  \label{lemma:balancing}
\end{lemma}

\paragraph{Mechanization.}
Mailbox configurations mostly rely on the definitions of list permutations
already provided by Rocq's standard library, but the formalization of mailbox
patterns requires more effort.
Describing the semantics of mailbox patterns as a relation between a configuration
and a pattern, instead of possibly infinite multisets, yields a mechanization that
is adequate to work with.  The proofs about pattern equivalence and inclusion
are typically inductions over a pattern, but proving balancing
(Lemma~\ref{lemma:balancing}) requires several technical lemmas that rely on
proven properties about permutations of mailbox configurations.

When mechanizing the pattern residual relation, we discovered that in the
definition of Pat~\cite{fowlerSpecialDeliveryProgramming2023} the rule for
$\mbstar{}$ was missing.  This issue was consequently fixed in the journal
version \cite{fowlerSpecialDeliveryProgramming2025}.

% Since the difference in mailbox types between Pat and the mailbox calculus are mostly the usage annotations provided by quasi-linear types,
% a mechanization of the mailbox calculus could certainly reuse our definitions and properties.
